
\subsubsection{2.1 Spurious Correlation and Shortcut Learning}

Geirhos et al. (2020) formalized the taxonomy of shortcut learning in deep neural networks, documenting how models exploit dataset biases rather than learning intended features. Sagawa et al. (2020) introduced Group DRO and the Waterbirds benchmark, enabling standardized evaluation of worst-group accuracy. Group DRO requires group annotations during training---a resource often unavailable in practice.

\subsubsection{2.2 Annotation-Free Robustification}

Liu et al. (2021) developed Just Train Twice (JTT), a two-stage approach that identifies spurious-reliant samples via training dynamics, then upweights them. JTT removes the need for group labels but requires two complete training runs. Kirichenko et al. (2023) showed that last-layer retraining (Deep Feature Reweighting, DFR) suffices for robustification, suggesting pretrained representations already contain core features. DFR requires held-out data with group labels for reweighting.

These methods operate on samples or final representations. None validates whether gradient directions themselves can distinguish spurious from core features---the assumption underlying gradient-based intervention.

\subsubsection{2.3 Simplicity Bias and Training Dynamics}

Shah et al. (2020) documented simplicity bias: neural networks trained with SGD preferentially learn simpler features first. Since spurious correlations often involve simpler patterns (background texture versus object shape), this predicts early spurious feature acquisition.

The present work attempts to measure this directly and finds the measurement apparatus itself requires validation.

\subsubsection{2.4 Gradient-Based Representation Analysis}

Parameter-space gradient subspace analysis appears primarily in optimization literature (incremental PCA, streaming SVD) rather than robustification. The present work applies subspace methods to spurious correlation measurement and identifies minimum requirements for such methods: sufficient accumulation density relative to parameter dimensionality.

